\section{Experiments and Results}
\label{sec:results}

\subsection{Per-row gradients are predominantly aligned on the rows that carry supervision}
\label{sec:aligned}
\paragraph{The honest denominator: active rows.} Because CE and the shared MTP
loss read the \emph{same} logits from the \emph{same} hidden state, a vocabulary
row that is not a target anywhere in the measurement batch has
$g^{\mathrm{mtp}}_i = 0.75\,g^{\mathrm{ce}}_i$ exactly (the two objectives differ
only in their target-indicator terms, and $0.75 = 0.5{+}0.25$ is the summed
auxiliary weight; the empirically measured norm ratio
$\lVert g^{\mathrm{mtp}}_i\rVert/\lVert g^{\mathrm{ce}}_i\rVert$ on such rows is
$0.750$, which is itself a check that all horizons use the same hidden state). Such rows
are scalar multiples with cosine identically $+1$, and they dominate the count. We
therefore classify rows by \emph{target presence}, not by a cosine threshold: a
row is \emph{active} iff $\lvert\lVert g^{\mathrm{mtp}}_i\rVert/\lVert
g^{\mathrm{ce}}_i\rVert - 0.75\rvert > 0.01$. By this construction
${\approx}88$--$92\%$ of rows are parallel and only \textbf{${\approx}8$--$12\%$
are active} (the exact fraction is batch-relative; see below). Restricted to the
active rows, the median per-row cosine is \textbf{${\approx}0.53$} (Muon
$0.529{\pm}0.006$, AdamW $0.547{\pm}0.013$, $n{=}3$), \textbf{${\approx}3\%$} are
opposed, and ${\approx}83\%$ exceed $\cos>0.3$. Over the full vocabulary the same
run reads a median of $+1.00$, $0.33\%$ opposed, and $98.4\%$ above $|\cos|>0.3$
against a $0.42\%$ calibrated control (Figure~\ref{fig:perrow}); these are the
construction-inflated numbers, reported only to calibrate against the control on a
common denominator. The active-row median of ${\approx}0.53$ is the honest measure
of directional alignment, and it is still well above zero, in sharp contrast to the
near-zero aggregate. Crucially, the active-row correction does \emph{not} affect
the mechanism in Section~\ref{sec:emergent}, which is mass-weighted: the parallel
rows carry only ${\approx}4\%$ of the gradient-magnitude mass, so the norm-profile
cosine ($0.98$) and opposed-norm fraction ($0.60$/$0.69$) are identical whether
computed over all rows or active rows only. The interference lives among the active rows;
restricting to them makes the finding sharper, not weaker.

\begin{figure}[t]
\centering
\includegraphics[width=\linewidth]{figures/fig1_perrow_histogram.pdf}
\caption{\textbf{Per-row CE/MTP gradients on the output projection are aligned on
the rows that carry supervision.} Main panel: distribution of \emph{signed} per-row
cosines over the active rows (median ${\approx}0.53$; ${\approx}3\%$ opposed, in
red). Inset: the full-vocabulary distribution, whose median is $+1.00$ because
${\approx}88$--$92\%$ of rows are parallel by construction ($g^{\mathrm{mtp}} =
0.75\,g^{\mathrm{ce}}$ for any row with no target in the batch, since CE and the
shared MTP loss read the same logits). Right: a CE-vs-L1 control reads $0.42\%$ of
rows above $|\cos|>0.3$ versus $98.4\%$ full-vocab for CE-vs-MTP; on active rows the
comparison is ${\approx}83\%$ versus the control, still a large calibration gap.
GPT-2 124M, step 1000, $50{,}257$ vocab rows (47 padding rows masked). The
mass-weighted mechanism (Section~\ref{sec:emergent}) is unaffected by the
active-row restriction because parallel rows carry only ${\approx}4\%$ of the
gradient mass.}
\label{fig:perrow}
\end{figure}

\subsection{The aggregate cosine decouples from per-row alignment during training}
\label{sec:emergent}
The aggregate cosine nonetheless reads $\approx 0$ ($-0.08$ at step 1000).
Tracking both metrics over training (Figure~\ref{fig:emergent}) shows this is not
a fixed property: at initialization the aggregate and per-row alignment
\textbf{agree at $\approx +0.99$}, and the aggregate then crosses zero early (between steps $100$ and $150$ under Muon, and between $150$ and $200$ under AdamW) while the per-row median holds at $\approx +1$ for the remainder of the
$1000$-step diagnostic window (late-training persistence is shown separately at
350M, Section~\ref{sec:rare}). These trajectories are a single diagnostic run per
optimizer; the close agreement between two different optimizers is the
replication we have.

We can pin down \emph{how} the aggregate collapses without any additional runs,
using an exact identity. The aggregate cosine is a magnitude-weighted average of
the per-row cosines, $\text{agg}=\sum_i w_i\cos_i$ with weights
$w_i=\lVert g^{\text{ce}}_i\rVert\lVert g^{\text{mtp}}_i\rVert/(\lVert
G^{\text{ce}}\rVert\lVert G^{\text{mtp}}\rVert)\ge 0$. If gradient magnitude were
spread evenly across rows, the weights would be equal and the aggregate would
equal the \emph{plain mean} of the per-row cosines. That mean is $+0.94$ (Muon)
and $+0.96$ (AdamW), whereas the actual aggregate is $-0.08$ and $-0.28$
(Figure~\ref{fig:normdecomp}). The entire $\approx\!1.0$ gap between ``aggregate
if norms were flat'' and the real aggregate is therefore attributable to the
per-row \emph{norm} structure, not to the typical per-row direction: the standard
diagnostic is reporting a magnitude-weighting artifact, which is the paper's
central point made quantitative.

The weights are unnormalized, and their sum is itself informative. Writing
$a_i=\lVert g^{\text{ce}}_i\rVert$ and $b_i=\lVert g^{\text{mtp}}_i\rVert$, we
have $\sum_i w_i=\langle a,b\rangle/(\lVert a\rVert\lVert b\rVert)=\cos(a,b)$,
which is exactly the \emph{norm-profile cosine} $\rho_{\text{norm}}$ defined
below. The aggregate therefore factors exactly as
\begin{equation}
\text{agg}=\sum_i w_i\cos_i
      =\underbrace{\Big(\textstyle\sum_i w_i\Big)}_{\rho_{\text{norm}}}
       \cdot\underbrace{\frac{\sum_i w_i\cos_i}{\sum_i w_i}}_{\langle\cos\rangle_{\text{mass}}}
      =\rho_{\text{norm}}\,\langle\cos\rangle_{\text{mass}},
\end{equation}
the product of how much the two losses' norm profiles overlap and the
mass-weighted mean of the per-row cosines. Numerically this is
$-0.08 = 0.98\times(-0.082)$ (Muon) and $-0.28 = 0.98\times(-0.286)$ (AdamW). The
two ``disambiguating statistics'' we introduce next are therefore not competing
hypotheses but the two exact factors of the aggregate: support overlap is the
total weight $\rho_{\text{norm}}$, and opposition is carried by
$\langle\cos\rangle_{\text{mass}}$.

What the identity does \emph{not} settle is which norm structure is responsible.
A negative aggregate requires the opposed minority ($\cos<0$, $\approx 0.33\%$ of
rows) to carry \emph{some} weight, so two mechanisms are a priori compatible with
the aligned-median-plus-negative-aggregate pattern: (a) \emph{support
divergence}, where CE and MTP load their magnitude onto increasingly disjoint
rows, making the products $w_i$ small precisely where each gradient is large; and
(b) a small number of \emph{high-norm opposed rows} dominating the weighted sum.
Separating (a) from (b) needs the per-row gradient \emph{norms}, which the earlier
runs do not store, so we measure them directly with a drop-in instrument
(\texttt{measure\_norm\_support.py}) that logs per-row norms and reports two
disambiguating statistics on the same short diagnostic pass: the
\emph{norm-profile cosine} (the cosine of the two per-row norm vectors; $\approx0$
under support divergence, $\approx1$ if the losses load the same rows) and the
\emph{opposed-norm fraction} (the share of $\lVert g^{\text{ce}}_i\rVert\lVert
g^{\text{mtp}}_i\rVert$ mass on rows with $\cos<0$). Across three seeds per
optimizer, the norm-profile cosine is $0.98$ (Muon $0.979\pm0.003$, AdamW
$0.977\pm0.007$): CE and MTP load the \emph{same} rows, so support does
\emph{not} diverge. Yet the opposed-norm fraction is $0.60\pm0.03$ (Muon) and
$0.69\pm0.01$ (AdamW): the majority of the magnitude mass sits on the opposed
minority. Mechanism (b) holds, and it holds under both optimizers, with AdamW
cancelling more strongly than Muon (Figure~\ref{fig:normsupport}); the difference
in opposed-norm fraction ($0.69$ vs $0.60$) is significant across seeds (Welch
$t{=}5.4$, $p\approx0.015$), so the magnitude is genuinely optimizer-dependent
while the mechanism is shared. The near-zero
aggregate is a norm-weighted cancellation by a high-norm opposed minority, not
disjoint support.

\begin{figure}[t]
\centering
\includegraphics[width=0.86\linewidth]{figures/fig7_norm_support.pdf}
\caption{\textbf{The aggregate is set by a high-norm opposed minority, not by
disjoint support} ($n{=}3$ seeds per optimizer, step $1000$, $50{,}257$ real
vocabulary rows). \emph{Left:} the norm-profile cosine (the cosine of the two
per-row gradient-norm vectors) is $\approx 0.98$ for both optimizers, so CE and
MTP load the same rows; support does not diverge. The actual aggregate cosine
nonetheless sits negative, and the arrow marks the collapse. \emph{Right:} the
opposed-norm fraction, the share of $\lVert g^{\text{ce}}_i\rVert\lVert
g^{\text{mtp}}_i\rVert$ mass on rows with $\cos<0$, is above one half for both
optimizers ($0.60$ Muon, $0.69$ AdamW): the opposed minority carries most of the
magnitude mass. AdamW cancels more strongly than Muon; the mechanism is identical,
the magnitude optimizer-dependent.}
\label{fig:normsupport}
\end{figure}

\paragraph{Which rows carry the cancellation.} The opposed mass is not spread
thinly: on the Muon seed-42 pass, $140$ opposed rows ($0.28\%$ of the vocabulary)
carry $60.7\%$ of the total $\lVert g^{\text{ce}}_i\rVert\lVert g^{\text{mtp}}_i\rVert$
mass, and just $47$ rows account for two-thirds of \emph{all} gradient mass
(Figure~\ref{fig:tokenlorenz}, left). Decoding those rows shows they are the most
frequent function words and punctuation: the four highest-mass opposed rows are
\texttt{.}, the newline, \texttt{\ the}, and \texttt{,}, which alone carry
${\approx}41\%$ of the mass, followed by \texttt{\ of}, \texttt{\ to}, \texttt{\ and},
and \texttt{\ a} (Figure~\ref{fig:tokenlorenz}, right). This is the expected locus:
these tokens are ubiquitous next-token \emph{and} future targets, and their
next-token and future-token continuation distributions genuinely differ, so the
CE and MTP gradients on their output rows point in different directions while both
being large. The cancellation is concentrated on a handful of high-frequency rows,
not diffuse across the tail.

\begin{figure}[t]
\centering
\includegraphics[width=\linewidth]{figures/fig_token_lorenz.pdf}
\caption{\textbf{The cancellation is carried by a few high-frequency rows.}
\emph{Left:} Lorenz curve of the per-row gradient mass
($\lVert g^{\text{ce}}_i\rVert\lVert g^{\text{mtp}}_i\rVert$, ranked descending);
$47$ rows account for two-thirds of all mass (dashed line), against the
line of equality (dotted). \emph{Right:} the highest-mass opposed rows, decoded,
are the most frequent function words and punctuation; the top four carry
${\approx}41\%$ of the total mass. Muon, seed 42, step 1000.}
\label{fig:tokenlorenz}
\end{figure}

\begin{figure}[t]
\centering
\includegraphics[width=0.72\linewidth]{figures/fig2_emergent_divergence.pdf}
\caption{\textbf{The aggregate and per-row alignment start together and
decouple.} Both metrics begin at $\approx +0.99$; within ${\sim}200$ steps the
aggregate (matrix-level) cosine falls below zero (crossing between steps
$100$--$150$ for Muon and $150$--$200$ for AdamW) while per-row alignment stays
high. The per-row curve is the full-vocabulary $|\cos|>0.3$ fraction
($\approx 0.98$); on active rows the level is lower ($\approx 0.83$, Section~\ref{sec:aligned})
but the decoupling from the aggregate is identical: the aggregate goes
negative under either denominator. One diagnostic run per optimizer (Muon solid,
AdamW dashed); the agreement between two optimizers is the replication.}
\label{fig:emergent}
\end{figure}

\begin{figure}[t]
\centering
\includegraphics[width=0.62\linewidth]{figures/fig6_norm_decomposition.pdf}
\caption{\textbf{Per-row norm structure, not the typical cosine, sets the
aggregate.} For each optimizer: the median per-row cosine ($+1.00$), the plain
\emph{mean} per-row cosine (which equals the aggregate the standard diagnostic
would report if gradient magnitude were spread evenly across rows), and the
\emph{actual} aggregate cosine. The mean is strongly positive ($+0.94$/$+0.96$)
while the actual aggregate is negative ($-0.08$/$-0.28$); the gap is due entirely
to per-row norm weighting. Single diagnostic run per optimizer, step $1000$,
$50{,}257$ real vocabulary rows.}
\label{fig:normdecomp}
\end{figure}

\subsection{The pattern is general, and survives untying}
\label{sec:general}
High per-row alignment is \textbf{not unique to the output head}. Across all 74
weight matrices (Figure~\ref{fig:perlayer}), attention and MLP matrices also show
high per-row alignment (many above 0.8), frequently with strongly \emph{negative}
aggregate cosine (down to $-0.78$). The output head is distinctive only in being
the point of \textbf{maximum} per-row alignment ($0.984$) sitting nearest
aggregate-zero, and, more importantly, in being the matrix that \emph{directly
produces the logits}. The head matters not because the metric singles it out, but
because interference there is paid directly in next-token loss. Alignment also
\textbf{persists when the head is untied} (A4: $99.3\%$ of rows above threshold,
versus $98.7\%$ tied at the matched step), so it is a property of the
output-projection gradient ($h^\top \partial\mathcal{L}/\partial\text{logits}$),
not an artifact of weight tying. These two percentages are \emph{full-vocabulary}
and therefore construction-floored (${\geq}75\%$ of rows are parallel by
construction); the untied run did not commit per-row norms, so we cannot restate
them on the active-row denominator, and we report them only as a point check that
untying does not \emph{reduce} alignment. The evidentiary weight of this check is
modest in any case: the soft-label identity (Section~\ref{sec:interventions}) predicts
untied-persistence \emph{a priori}, since that identity never references tying.
The per-layer generality (Figure~\ref{fig:perlayer})
and the untied-head comparison are each \emph{single-run} measurements at a fixed
step; we report them as point observations, not seed-replicated estimates (only
the 124M CE-vs-MTP loss gap is replicated across seeds).

We are deliberate about what this generality does and does not license. The
same qualitative signature (high per-row alignment with a non-positive aggregate) appears across the network, and at many matrices the aggregate is
\emph{strongly} negative (to $-0.78$), not merely near zero. This means the
per-row/aggregate gap is a fairly generic property of these structured gradient
matrices, so a near-zero aggregate on its own is \emph{not} a specific signature
of two-loss interference. What is specific to the output head is that it is the
one matrix whose per-row gradients are the CE and MTP gradients directly, whose
norm-weighted opposition we measure end to end (Figure~\ref{fig:normsupport}),
and whose interference is paid immediately in next-token loss. We do not
adjudicate the per-row norm structure of the strongly-negative aggregates
elsewhere in the network; a per-row account of those is beyond what our
measurements settle.

\begin{figure}[t]
\centering
\includegraphics[width=0.66\linewidth]{figures/fig3_perlayer.pdf}
\caption{\textbf{High per-row alignment is general across the network.} Per-row
alignment vs aggregate cosine for all 74 weight matrices. The output head (star)
is the maximum and sits nearest aggregate-zero, but attention/MLP matrices show
the same pattern, many with strongly negative aggregate cosine.}
\label{fig:perlayer}
\end{figure}

\subsection{Adding shared MTP degrades next-token loss; no auxiliary recovers it}
\label{sec:interventions}
Adding the shared MTP objective degrades next-token validation loss from
$\mathbf{3.985 \pm 0.014}$ (CE-only, $n{=}5$) to $\mathbf{4.377 \pm 0.021}$
($n{=}3$): a \textbf{$+0.39$ nat} (${\approx}10\%$) increase. On the three seeds
the two conditions share ($42,123,456$), a paired test gives a mean gap of
$\mathbf{0.382}$ nats (per-seed $0.397/0.380/0.370$), $t{=}50.0$,
$\mathbf{p \approx 4.0\times10^{-4}}$ ($\text{df}{=}2$); treating the conditions
as independent (Welch, CE-only minus shared-MTP, using all $n{=}5$ and $n{=}3$
seeds) gives $t = -28.8$ (negative because CE-only is the lower loss),
$p \approx 8\times10^{-5}$ with a $95\%$ CI on the difference of
$[0.35, 0.43]$ nats. The gap is large relative to the seed-to-seed spread
(within-condition sd $\approx 0.01$--$0.02$), which is what makes it detectable
at these small $n$; we report the effect as the $0.39$-nat loss increase itself
rather than a standardized effect size, since the latter is inflated by the tight
reproducibility of a controlled testbed. With $n{=}3$ per condition and this
spread, the minimum difference we have ${\approx}80\%$ power to detect is
$\Delta\approx0.05$ nats; the headline gap is ${\sim}8\times$ that floor, so it is
detected with room to spare, whereas differences below ${\approx}0.05$ nats
(relevant to the tuned auxiliary below) are at or under our detection limit. We
report uncorrected $p$-values for the $k{=}5$ variant-vs-baseline comparisons;
under a Holm correction at $\alpha{=}0.05$ the shared-MTP, stop-gradient,
standard-contrastive, and NextLat results remain significant, while only the
tuned-auxiliary comparison ($p\approx0.19$) does not survive, consistent with its
effect sitting below the detection floor. Detaching
the auxiliary from the head via stop-gradient makes it \emph{worse} ($4.493$,
$\Delta{+}0.508$, $p < 10^{-4}$), not better, which points to a norm-weighting/capacity account rather than directional conflict (we note this is a data
point, not a proof: stop-gradient also removes the auxiliary's ability to shape
the shared representation, which is a competing reason it could hurt). The
$n{=}5$ CE-only and $n{=}3$ shared-MTP conditions share three seeds ($42,123,456$);
the two extra CE-only seeds ($789,1337$) were run for a tighter baseline and no
MTP seeds were dropped. Capacity-separated auxiliaries reduce but do \textbf{not}
recover the baseline. The \emph{standard} layer-split contrastive auxiliary
($\mathrm{G}_{\text{nce}}$: an InfoNCE-style contrastive loss applied to an
intermediate layer, giving the auxiliary its own representation capacity instead
of sharing the output head) reaches $4.024$ ($\Delta{+}0.039$, $p \approx 0.010$),
statistically distinguishable from A, i.e.\ a genuine non-recovery. A
hyperparameter-tuned variant reaches $4.003$ ($\Delta{+}0.018$) but is \emph{not}
statistically separable from A at $n{=}3$ ($p \approx 0.19$): for this variant we
can claim neither recovery nor non-recovery, and we rest the capacity conclusion
on the standard variant. A latent-prediction baseline (NextLat: the auxiliary
predicts the \emph{hidden state} at future offsets rather than their tokens, so it
never touches the output head) lands at $4.037$ ($\Delta{+}0.051$, Welch
$p{\approx}0.018$) (Figure~\ref{fig:interventions}).

\begin{figure}[t]
\centering
\includegraphics[width=0.66\linewidth]{figures/fig4_interventions.pdf}
\caption{\textbf{No auxiliary recovers the baseline; MTP degrades it most.} Final
next-token validation loss by variant (points = seeds). Shared MTP degrades CE by
$\Delta{+}0.39$; stop-grad by $\Delta{+}0.51$; no capacity-separated auxiliary
recovers the CE-only baseline within noise.}
\label{fig:interventions}
\end{figure}

Conflict-targeted gradient surgery also does not recover the baseline, and this
is a stronger statement than ``there was nothing to remove.'' The opposed rows
are negligible by count ($0.3\%$) but carry the majority of the gradient mass
($60$--$69\%$), and a projection acts on the mass-weighted vector, so surgery had
a great deal to project away; per-row Gram--Schmidt on the head targets exactly
those opposed rows and still recovers nothing. The measured opposition is
mass-dominant and causally inert. Because those runs are single-seed with a
numerical-drift caveat, we report them in full in Appendix~\ref{app:surgery}
rather than resting any claim on them.

\subsection{The apparent rare-token alignment is a construction artifact; the pattern holds at 350M}
\label{sec:rare}
Measured over the full vocabulary, per-row alignment appears \emph{higher} for
rarer tokens: the rarest quintile averages $0.98$ mean $|\cos|$ versus $0.86$
for the most-common quintile (Figure~\ref{fig:tokenfreq}, left). This effect is
an \emph{artifact of the construction floor}. Rare tokens are almost never
targets in a measurement batch, so they are overwhelmingly parallel-by-construction
($\cos=+1$): the active fraction falls from ${\approx}25\%$ in the common quintile
to ${\approx}1.5\%$ in the rarest. Restricted to active rows, the apparent
gradient disappears: active-row mean $|\cos|$ is essentially flat across
frequency quintiles ($0.51$--$0.54$ for Muon, $0.52$--$0.60$ for AdamW;
Figure~\ref{fig:tokenfreq}, right), with only a weak residual (Spearman
$\rho{\approx}0.12$ between token id and active-row $|\cos|$). We therefore
\emph{retract} the rare-token-alignment reading: the full-vocabulary ordering
reflects how many rows are parallel by construction, not a frequency-dependent
alignment. The
aligned-per-row-with-negative-aggregate picture holds qualitatively at
\textbf{350M} (single seed): CE-only $3.752$, shared-MTP $4.129$ ($\Delta{+}0.378$).

The 350M runs are the paper's only \emph{late-training} per-row measurement,
logged out to step $87{,}000$. The instrument computes the same quantity in both
runs (the per-row cosine between the CE gradient and the MTP gradient on the output projection) regardless of which objective the model was trained with
(the MTP gradient is computed at measurement time for both). Measured this way,
the shared-MTP-trained model holds $|\cos|>0.3$ at $98.9\%$ of rows at step
$87{,}000$, and the CE-only-trained model reads $94.5\%$ at the same step. These
are \emph{full-vocabulary} figures and thus construction-floored: with ${\geq}75\%$
of rows parallel by construction, the informative range is the remaining
${\approx}25$ percentage points, within which the MTP-trained model sits near the
top and the CE-only model ${\approx}4$ points lower (coarse, but not vacuous),
stated here at its true resolution (the 350M runs committed only summary statistics,
so an active-row restatement is not recoverable from the saved data). That
the CE-only model \emph{also} shows high CE/MTP alignment this late indicates the
alignment is an intrinsic property of the output-projection gradient geometry,
not something the MTP objective induces. Both are single seeds, so we treat the
per-row persistence as illustrative (consistent with alignment not decaying over a
long horizon), not established. The 350M \emph{loss} gap ($\Delta{+}0.378$ above)
is a next-token-loss number, unaffected by the active-row denominator, and it is
this scale-stability of the gap (${+}0.392$ at 124M, ${+}0.378$ at 350M) that
Section~\ref{sec:interventions} builds on.

\begin{figure}[t]
\centering
\includegraphics[width=0.6\linewidth]{figures/fig5_token_frequency.pdf}
\caption{\textbf{The rare-token alignment effect is a construction artifact.}
Left: over the full vocabulary, mean per-row $|\cos|$ rises from common to rare
tokens (both optimizers); this is the apparent effect. Right: restricted to active
rows (those carrying a target), the effect vanishes and mean $|\cos|$ is flat
across frequency quintiles; the grey bars show the active fraction collapsing
from ${\approx}25\%$ (common) to ${\approx}1.5\%$ (rare), which is what drives the
full-vocabulary trend. Single run per optimizer, step 1000, GPT-2 124M.}
\label{fig:tokenfreq}
\end{figure}
